\documentclass[11pt]{article}

\usepackage[margin=1in]{geometry}
\usepackage{amsmath,amssymb,amsthm}
\usepackage[hidelinks]{hyperref}

% Theoretical-node template. Replace every placeholder with valid LaTeX and
% keep only the theorem environments used by the instantiated node.
\newtheorem{theorem}{Theorem}
\newtheorem{lemma}[theorem]{Lemma}
\newtheorem{proposition}[theorem]{Proposition}
\newtheorem{corollary}[theorem]{Corollary}

\theoremstyle{definition}
\newtheorem{definition}[theorem]{Definition}
\newtheorem{assumption}{Assumption}

\title{[{{NODE_ID}}] {{NODE_TITLE}}}
\author{}
\date{}

\begin{document}
\maketitle

\section*{Setup}
% Define every object and piece of notation needed to understand the statement.
% Use definition environments for substantive mathematical definitions.
{{SETUP_BODY}}

\section*{Assumptions}
% Give each assumption a precise, referenceable environment. Repeat this block
% as needed and make the environments agree one-for-one with Nxxx.yaml.
\begin{assumption}[{{ASSUMPTION_NAME}}]\label{ass:{{NODE_ID}}-1}
{{ASSUMPTION_BODY}}
\end{assumption}

\section*{Main Statement}
% Use exactly one primary theorem-like statement for this node. Replace theorem
% by lemma, proposition, or corollary when that better reflects its mathematical
% role; do not duplicate this presentation choice in Nxxx.yaml.
% Quantify every object, invoke every required assumption explicitly, and state
% the exact conclusion without relying on the originating task artifact.
\begin{theorem}[{{STATEMENT_NAME}}]\label{stmt:{{NODE_ID}}}
{{FORMAL_STATEMENT}}
\end{theorem}

\paragraph{Interpretation.}
{{INTERPRETATION}}

\paragraph{Proof idea.}
{{PROOF_IDEA}}

% For a citation route, use the proof environment below to state the exact
% cited theorem and verify each hypothesis and quantifier before applying it.
\section*{Proof}
\begin{proof}
{{COMPLETE_PROOF}}
\end{proof}

% With several alternative proofs, replace the single Proof idea and Proof
% block by one descriptively titled section per proof. Put an accurate Proof
% idea paragraph immediately before each proof environment, and make each
% section title match the corresponding supports[].tex_section in Nxxx.yaml.
% Whenever a proof invokes a prerequisite node, state the exact result used and
% cite its N-ID in the mathematical prose; the YAML dependency list is not a
% substitute for a mathematically explicit invocation. Do not render DAG
% terms such as support-route kind, dependencies, audit paths, or provenance in
% this mathematical document; keep them in YAML. State and prove every needed
% intermediate claim, justify every inference, cover all relevant boundary and
% degenerate cases, and never hide a gap behind "obvious", "standard",
% "routine", or "left to the reader". Split long arguments into named lemmas
% with complete proofs. The interpretation and proof idea must be correct and
% readable, but neither may carry any logical step omitted from the proof.

\end{document}
